\documentclass[11pt,a4paper]{article}

\usepackage[margin=2.5cm]{geometry}
\usepackage[brazilian]{babel}
\usepackage[T1]{fontenc}
\usepackage[utf8]{inputenc}
\usepackage{microtype}
\usepackage{amsmath}
\usepackage{booktabs}
\usepackage{longtable}
\usepackage{tabularx}
\usepackage{enumitem}
\usepackage{xcolor}
\usepackage{csquotes}
\usepackage{hyperref}
\usepackage[backend=biber,style=authoryear,maxcitenames=2,maxbibnames=99,
  doi=true,isbn=false,url=false]{biblatex}

\addbibresource{mpga_rsl.bib}

\hypersetup{
  colorlinks=true,
  linkcolor=blue!50!black,
  citecolor=blue!50!black,
  urlcolor=blue!50!black,
  pdftitle={Protocolo de revisão sistemática e análise bibliométrica sobre métodos de medição de material particulado}
}

\setlength{\parindent}{0pt}
\setlength{\parskip}{0.65\baselineskip}
\setlist[itemize]{topsep=2pt,itemsep=2pt,parsep=0pt}
\setlist[enumerate]{topsep=2pt,itemsep=2pt,parsep=0pt}
\renewcommand{\arraystretch}{1.18}

\newcommand{\MPdez}{MP$_{10}$}
\newcommand{\MPdoiscinco}{MP$_{2,5}$}
\newcommand{\versao}{0.5}
\newcommand{\dataversao}{5 de outubro de 2026}
\newcommand{\pendente}[1]{%
  \begin{center}
  \fcolorbox{orange!70!black}{orange!8}{%
    \parbox{0.92\textwidth}{\textbf{Decisão a ratificar antes do registro:} #1}}
  \end{center}}

\title{\textbf{Protocolo de revisão sistemática}\\[0.6em]
\large Comparabilidade entre métodos automáticos e gravimétricos na medição de
\MPdez{} e \MPdoiscinco{} no ar ambiente: revisão sistemática, análise
bibliométrica descritiva e síntese dos fatores condicionantes}
\author{André Henrique Oliveira Santos\\
\small Revisão conduzida por um único pesquisador}
\date{Versão \versao{} --- \dataversao}

\begin{document}

\maketitle

\begin{center}
\fcolorbox{red!60!black}{red!5}{%
  \parbox{0.92\textwidth}{\textbf{Situação do documento.} Protocolo não registrado,
  emendado após a busca-piloto e após a seleção, a extração inicial e a avaliação
  metodológica piloto conduzidas por um pesquisador assistido por IA. As buscas em
  Scopus e Web of Science foram executadas
  em 26 de setembro de 2026, e a busca na PubMed em 5 de outubro de 2026. O instrumento
  metodológico já havia sido aplicado em caráter piloto aos 40 documentos incluídos
  quando a Emenda 3 foi elaborada; 17 publicações originais haviam sido classificadas
  como definitivamente não recuperadas. A revisão seguirá com pesquisador único e
  incluirá análise bibliométrica descritiva. Portanto, um registro realizado a partir
  desta versão não será apresentado como registro prospectivo. Todas as decisões tomadas após o
  conhecimento dos resultados deverão permanecer datadas e identificadas como
  emendas.}}
\end{center}

\section*{Controle de versões}

\begin{tabularx}{\textwidth}{@{}p{1.5cm}p{3.2cm}X@{}}
\toprule
\textbf{Versão} & \textbf{Data} & \textbf{Descrição} \\
\midrule
0.1 & 26/09/2026 & Estrutura inicial do protocolo, pergunta, critérios de
elegibilidade, plano de busca, extração, avaliação metodológica e síntese. \\
0.2 & 27/09/2026 & Emenda após a busca-piloto e a primeira avaliação dos estudos:
registro transparente do estágio da revisão, buscas adicionais, uso de IA, segunda
revisão independente e regras do instrumento metodológico. \\
0.3 & 05/10/2026 & Execução da busca na PubMed e delimitação das bases bibliográficas
a Scopus, Web of Science e PubMed. SciELO, Compendex e Embase foram retiradas do plano
antes da execução, com manutenção das buscas retrospectiva e prospectiva de citações. \\
0.4 & 05/10/2026 & Atualização do andamento após a avaliação integral dos seis
relatórios novos da PubMed: cinco inclusões e uma exclusão, sem alteração dos
critérios de elegibilidade ou do plano de síntese. \\
0.5 & 05/10/2026 & Emenda metodológica consolidada: revisão por pesquisador único,
segunda conferência interna, encerramento da recuperação de 17 publicações, inclusão
de bibliometria descritiva sem filtragem multicritério e revisão das unidades e regras
da avaliação metodológica e do tratamento de estimativas dependentes. \\
\bottomrule
\end{tabularx}

\section{Informações administrativas}

\subsection{Título da revisão}

Comparabilidade entre métodos automáticos e gravimétricos na medição de \MPdez{}
e \MPdoiscinco{} no ar ambiente: revisão sistemática, análise bibliométrica
descritiva e síntese dos fatores condicionantes.

\subsection{Equipe, responsabilidades e contato}

A revisão será conduzida por André Henrique Oliveira Santos, responsável pela busca,
seleção, recuperação dos documentos disponíveis, extração, avaliação metodológica,
análise bibliométrica, síntese e redação. Não haverá segundo revisor humano
independente. As etapas sujeitas a julgamento serão conferidas pelo próprio
pesquisador em uma segunda rodada separada, com registro das alterações efetuadas.

Eventuais contribuições adicionais serão descritas no manuscrito final segundo a
taxonomia CRediT e somente resultarão em autoria quando os critérios editoriais de
autoria forem atendidos.

\subsection{Uso de ferramentas de inteligência artificial}

O Codex, da OpenAI, foi utilizado a partir de 26 de setembro de 2026 para apoiar a
organização dos arquivos RIS, a deduplicação, a preparação de planilhas, a triagem
preliminar, a extração inicial de dados, a leitura assistida de figuras e a avaliação
metodológica piloto, além da preparação de rotinas reprodutíveis de análise. A
ferramenta não será considerada autora, revisora independente ou fonte de decisão
metodológica. A responsabilidade pelas decisões de elegibilidade, pelos valores
extraídos, pelos julgamentos metodológicos e pela síntese será do pesquisador. Os
resultados assistidos por IA serão conferidos diretamente contra as fontes em uma
segunda rodada interna, e as correções serão registradas.

\subsection{Registro e acesso ao protocolo}

O protocolo será depositado no OSF Registries, preferencialmente utilizando o modelo
\emph{Generalized Systematic Review}. O registro informará que, antes do depósito,
foram realizadas as buscas em Scopus, Web of Science e PubMed, a primeira triagem, a
avaliação de textos completos, a extração preliminar e a avaliação metodológica
piloto. O
endereço e o identificador serão inseridos nesta seção e no resumo do artigo. O
depósito não será descrito como pré-registro ou registro prospectivo.

\subsection{Emendas}

Cada emenda deverá informar data, seção afetada, descrição, justificativa, momento
da revisão em que ocorreu e possível efeito sobre os resultados. Mudanças realizadas
após o conhecimento dos resultados não poderão ser apresentadas como decisões
definidas previamente.

A primeira emenda foi elaborada em 27 de setembro de 2026, após a recuperação e a
avaliação preliminar dos resultados de Scopus e Web of Science. Ela documenta o estágio
real da revisão, determina a execução de buscas adicionais e exige uma segunda revisão
independente. A segunda emenda, de 5 de outubro de 2026, registra a execução da PubMed
e a delimitação das bases bibliográficas. A terceira emenda, também de 5 de outubro de
2026, substitui a exigência de duplicação humana por uma segunda conferência interna,
encerra a tentativa de recuperação de 17 publicações, incorpora análise bibliométrica
descritiva sem seleção por indicadores de impacto e consolida as unidades e as regras
da avaliação metodológica e do tratamento de estimativas dependentes após o
teste-piloto. A Emenda 3 prevalece sobre a exigência de segundo revisor constante da
Emenda 1. Os registros completos das emendas serão disponibilizados com o protocolo.

\subsection{Apoio, financiamento e conflitos de interesse}

As fontes de apoio financeiro e não financeiro, a participação dos financiadores e
os conflitos de interesse dos autores serão declarados antes do registro. Não haverá
participação de fabricantes de equipamentos nas decisões de elegibilidade, extração,
avaliação ou síntese.

\section{Justificativa}

Os métodos gravimétricos determinam a concentração de material particulado a partir
da massa coletada em filtro e do volume de ar amostrado durante um período definido.
Os métodos automáticos produzem observações em maior resolução temporal por meio de
princípios como atenuação beta, microbalança oscilante e detecção óptica. As duas
categorias de método podem diferir em razão do princípio de medição, do sistema de
entrada, do condicionamento da amostra, da perda de componentes semivoláteis, da
umidade, da composição do aerossol e dos procedimentos de calibração e correção.

Estudos de campo frequentemente avaliam equipamentos automáticos por comparação
simultânea com amostradores gravimétricos. Entretanto, os resultados podem ser
expressos como diferenças médias, vieses relativos, razões, regressões, correlações,
limites de concordância ou atendimento a critérios regulatórios. Essa diversidade
dificulta a interpretação conjunta da comparabilidade e a identificação das condições
em que os métodos podem ou não produzir resultados intercambiáveis.

A revisão pretende organizar essa evidência sem antecipar a pergunta empírica da
futura dissertação. Enquanto a dissertação avaliará a incerteza decorrente da escolha
temporal dos dias de amostragem, esta revisão examinará a incerteza associada aos
próprios métodos de medição.

O relatório final seguirá o PRISMA 2020 e seu documento de explicação e elaboração
\parencite{page_prisma_statement_2021,page_prisma_explanation_2021}. Este protocolo
será conferido com o PRISMA-P, e as estratégias de busca serão documentadas segundo
o PRISMA-S. Caso a combinação quantitativa não seja apropriada, a apresentação da
síntese seguirá também os princípios do SWiM. As CEE Guidelines serão utilizadas como
fonte de orientação, mas não será reivindicada conformidade integral, pois o protocolo
não foi registrado antes do início da revisão e não haverá avaliação humana
independente. A análise bibliométrica terá finalidade descritiva e não será usada para
determinar elegibilidade, qualidade ou peso dos estudos.

\section{Objetivos}

\subsection{Objetivo geral}

Avaliar sistematicamente a comparabilidade entre métodos automáticos e gravimétricos
na determinação de concentrações de \MPdez{} e \MPdoiscinco{} no ar ambiente,
quantificando os vieses relatados e identificando fatores ambientais, operacionais e
instrumentais associados às diferenças entre os métodos.

\subsection{Objetivos específicos}

\begin{enumerate}[label=\alph*)]
  \item caracterizar os locais, períodos, equipamentos, métodos de referência e
  desenhos dos estudos comparativos;
  \item reunir e padronizar, quando possível, medidas de diferença, viés, razão,
  concordância, regressão, erro e equivalência;
  \item comparar separadamente os resultados referentes a \MPdez{} e
  \MPdoiscinco{};
  \item avaliar diferenças entre tecnologias automáticas, incluindo BAM, TEOM,
  TEOM-FDMS e outros métodos elegíveis;
  \item investigar a influência de umidade, temperatura, sazonalidade, composição
  do aerossol, faixa de concentração, tipo de local, tratamento da amostra e
  correções instrumentais;
  \item caracterizar a evolução temporal, as fontes, os países, os autores, as
  palavras-chave e as tecnologias presentes no conjunto de documentos incluídos;
  \item avaliar as limitações metodológicas e a consistência da evidência disponível;
  \item identificar lacunas relevantes para o monitoramento ambiental e para estudos
  futuros de representatividade temporal.
\end{enumerate}

\section{Pergunta da revisão e estrutura conceitual}

\textbf{Pergunta:} Em estudos de monitoramento do ar ambiente com equipamentos
operados simultaneamente e no mesmo local, qual é a concordância entre as
concentrações de \MPdez{} e \MPdoiscinco{} determinadas por métodos automáticos e
por métodos gravimétricos de referência, e quais fatores ambientais, operacionais e
instrumentais influenciam as diferenças observadas?

Será utilizada uma estrutura PICOS adaptada a estudos de comparação de métodos.

\begin{longtable}{@{}p{3.6cm}p{10.9cm}@{}}
\toprule
\textbf{Componente} & \textbf{Definição} \\
\midrule
\endhead
População/contexto & Locais de monitoramento de ar ambiente externo nos quais foram
realizadas medições de \MPdez{} e/ou \MPdoiscinco{}. \\
Método índice & Equipamento automático ou contínuo destinado à determinação da
concentração em massa de material particulado. \\
Comparador & Método com coleta em filtro e determinação gravimétrica da massa,
descrito como método de referência, equivalente ou comparador validado. \\
Desfechos & Diferença, viés, razão, concordância, erro, regressão, precisão,
correlação e atendimento a critérios de equivalência. \\
Desenho & Estudo primário de comparação em campo, com medições simultâneas ou
temporalmente coincidentes e equipamentos colocalizados ou representativos da mesma
massa de ar. \\
\bottomrule
\end{longtable}

\section{Critérios de elegibilidade}

\subsection{Critérios de inclusão}

Serão elegíveis estudos que atendam conjuntamente aos seguintes critérios:

\begin{enumerate}
  \item apresentem dados primários obtidos em monitoramento de ar ambiente externo;
  \item avaliem \MPdez{}, \MPdoiscinco{} ou ambas as frações;
  \item comparem um método automático com um método gravimétrico baseado em filtro;
  \item realizem as medições em períodos coincidentes e em equipamentos colocalizados
  ou posicionados de modo que representem justificadamente a mesma massa de ar;
  \item informem dados quantitativos de comparabilidade ou dados suficientes para
  calcular ao menos uma medida de comparação;
  \item descrevam suficientemente os equipamentos e os períodos de amostragem para
  permitir a interpretação dos resultados;
  \item sejam artigos completos revisados por pares ou relatórios técnicos oficiais
  com método e resultados primários suficientemente documentados.
\end{enumerate}

Não será estabelecida duração mínima a priori. A duração, o número de pares, a
completude e a cobertura sazonal serão extraídos e considerados na avaliação
metodológica e nas análises de sensibilidade.

\subsection{Critérios de exclusão}

Serão excluídos:

\begin{itemize}
  \item estudos exclusivamente laboratoriais, de câmara ou de bancada;
  \item monitoramento interno, ocupacional, pessoal ou de emissões em chaminés;
  \item estudos que comparem somente dois equipamentos automáticos ou somente dois
  métodos gravimétricos;
  \item estudos sem coincidência temporal entre os métodos;
  \item estudos de composição química sem comparação da concentração em massa;
  \item estudos de partículas totais em suspensão sem resultados separados para
  \MPdez{} ou \MPdoiscinco{};
  \item resumos de eventos, editoriais, cartas sem dados primários, revisões e
  protocolos;
  \item publicações duplicadas do mesmo estudo sem resultados adicionais.
\end{itemize}

Sensores classificados como de baixo custo serão excluídos da análise principal. Um
equipamento compacto somente será elegível quando o estudo o apresentar como método
de monitoramento ambiental de desempenho equivalente, de referência ou
\emph{research grade}, com documentação de calibração e comparação direta com o
método gravimétrico.

\pendente{Ratificar a inclusão de relatórios técnicos oficiais como evidência
primária. A recomendação inicial é incluí-los em um estrato separado dos artigos
revisados por pares, reduzindo o risco de perda de estudos de equivalência sem misturar
automaticamente fontes com processos editoriais diferentes.}

\subsection{Idiomas, período e situação de publicação}

Não serão aplicados filtros de idioma ou data durante as buscas. Textos em idiomas
não dominados pelo pesquisador serão avaliados com tradução assistida e, quando necessário,
apoio externo. A data final corresponderá ao dia de execução de cada busca. A busca
será atualizada próximo à submissão do manuscrito.

\section{Fontes de informação}

As seguintes bases integram o plano emendado:

\begin{itemize}
  \item Scopus, busca executada em 26 de setembro de 2026;
  \item Web of Science Core Collection, busca executada em 26 de setembro de 2026;
  \item MEDLINE via PubMed, busca executada em 5 de outubro de 2026.
\end{itemize}

Na PubMed foram recuperados 64 registros sem filtros de idioma, data ou tipo
documental. A deduplicação confirmou 31 coincidências por DOI e duas por título,
autores, ano e periódico. As duas coincidências sem DOI foram confirmadas pelo usuário
em 5 de outubro de 2026. Assim, 33 duplicatas foram removidas e 31 registros foram
encaminhados para triagem. A triagem de título e resumo excluiu 25 registros e
encaminhou seis para localização e avaliação do texto completo.

Os seis relatórios foram recuperados e avaliados integralmente em 5 de outubro de
2026. Foram incluídos RSL-126, RSL-129, RSL-136, RSL-144 e RSL-145. O RSL-143 foi
excluído porque o QT-50 é apresentado como sensor óptico de baixo custo, sem
classificação como método ambiental equivalente, de referência ou
\emph{research grade}, conforme o critério específico estabelecido neste protocolo.

As estratégias preparadas para SciELO, Compendex e Embase na versão 0.5 não foram
executadas e foram retiradas do plano corrente em 5 de outubro de 2026. Essa decisão
ocorreu após o início da revisão e será declarada como emenda e como possível limitação
de cobertura no relatório final.

Considerando as três bases, foram identificados 233 registros brutos e removidas 83
duplicatas. Dos 150 registros únicos triados, 90 foram excluídos por título e resumo e
60 publicações originais foram procuradas para leitura integral. Dezessete publicações
originais não puderam ser obtidas e as tentativas de recuperação foram encerradas em
5 de outubro de 2026. Esses registros serão apresentados como relatórios não
recuperados, sem atribuição de motivo de exclusão por conteúdo. Foram avaliadas 43
publicações originais e uma fonte alternativa; quatro documentos foram excluídos e 40
foram incluídos, representando no máximo 39 estudos independentes antes das buscas de
citações.

Serão examinadas as listas de referências dos estudos incluídos e das revisões
relevantes. Também será realizada busca de citações posteriores dos estudos incluídos.
Caso os relatórios oficiais permaneçam elegíveis, serão definidos previamente os
portais institucionais e os procedimentos reproduzíveis de busca.

Autores poderão ser contatados quando informações essenciais de método ou resultados
dos documentos incluídos não puderem ser obtidas no texto e nos suplementos. Serão
registradas a data, a forma de contato e a resposta recebida. Não serão realizadas
novas tentativas de obtenção das 17 publicações classificadas nesta emenda como não
recuperadas.

\section{Estratégia de busca}

A estratégia combinará cinco blocos conceituais:

\begin{enumerate}
  \item material particulado e suas frações;
  \item método gravimétrico ou baseado em filtro;
  \item método automático ou tecnologia instrumental;
  \item comparação, concordância, viés ou equivalência;
  \item ar ambiente e qualidade do ar.
\end{enumerate}

A versão 0.3 efetivamente executada na Scopus é apresentada a seguir. A mesma estrutura
foi adaptada à Web of Science e, com termos livres e descritores MeSH, à PubMed. Os
conceitos essenciais e cada comando completo foram registrados separadamente:

\begin{verbatim}
TITLE-ABS-KEY (
  ("particulate matter" OR PM10 OR "PM 10"
   OR PM2.5 OR "PM 2.5" OR PM25)
  AND
  (gravimetric OR "filter-based" OR "filter based"
   OR "reference method" OR "Federal Reference Method" OR FRM
   OR "manual sampler" OR "filter sampler"
   OR "high volume sampler" OR "low volume sampler")
  AND
  ("automatic monitor" OR "automatic monitoring"
   OR "continuous monitor" OR "continuous monitoring"
   OR "real-time monitor" OR BAM OR "beta attenuation monitor"
   OR TEOM OR "tapered element oscillating microbalance"
   OR FDMS OR "filter dynamics measurement system"
   OR SHARP OR nephelometer OR "optical monitor")
  AND
  (compare OR comparison OR agreement OR concordance OR equivalence OR bias
   OR validation OR evaluation OR intercomparison OR collocated
   OR "co-located" OR "side-by-side" OR performance)
  AND
  ("ambient air" OR "air quality")
)
\end{verbatim}

A busca-piloto foi realizada em Scopus e Web of Science sem artigos-semente, e a busca
adicional foi executada na PubMed. A calibração utilizou quantidade e relevância dos
resultados, sinônimos encontrados e motivos de irrelevância. A recuperação dos estudos
incluídos será verificada retrospectivamente como avaliação de cobertura, sem apresentar
esse procedimento como validação prospectiva. As estratégias não receberão revisão
independente por segundo pesquisador; essa ausência será declarada como limitação.

PTS/TSP poderá ser acrescentada apenas a uma busca-piloto exploratória. Sua inclusão
na revisão principal exigirá nova emenda anterior à incorporação dos resultados e
apresentação separada das frações \MPdez{} e \MPdoiscinco{}.

Para cada fonte serão arquivados nome da base e plataforma, comando completo, data,
limites, número de registros e formato de exportação.

\section{Gerenciamento dos registros e seleção}

O arquivo \texttt{mpga\_rsl.bib} contém referências metodológicas e modelos de artigo
gerenciados no Zotero. Ele é somente leitura neste projeto e não será editado,
deduplicado ou ampliado automaticamente. Os resultados das buscas da revisão serão
mantidos em coleção e exportações próprias, separadas desse arquivo.

Após a reunião dos resultados, duplicatas serão identificadas por DOI, título, autores,
ano e conferência manual. O número removido antes da triagem será registrado no
fluxograma PRISMA.

A seleção ocorrerá em duas etapas, conduzidas pelo pesquisador responsável:

\begin{enumerate}
  \item leitura de títulos e resumos;
  \item leitura dos textos completos considerados potencialmente elegíveis.
\end{enumerate}

A seleção já realizada será considerada a primeira rodada de um pesquisador assistido
por IA. O mesmo pesquisador realizará uma segunda conferência dos registros em rodada
separada e, quando operacionalmente possível, sem visualizar inicialmente a decisão
anterior. Mudanças entre as rodadas serão resolvidas pela releitura do registro ou do
texto completo e acompanhadas de justificativa. Na etapa de texto completo, será
atribuído um único motivo principal de exclusão, segundo uma lista padronizada. Não
serão calculadas medidas de concordância entre revisores, pois não haverá avaliações
humanas independentes.

As 17 publicações não recuperadas permanecerão separadas das exclusões. Não será
inferida sua elegibilidade a partir do título ou do resumo, e seu possível efeito sobre
a completude da revisão será discutido como limitação.

Relatórios múltiplos de um mesmo estudo serão vinculados a uma identificação comum.
O estudo, e não cada publicação isolada, será a unidade principal da síntese.

Será adotada a seguinte hierarquia de identificação: família de estudo, relatório,
comparação elegível e estimativa. Uma comparação elegível corresponde a uma combinação
específica de local, período, fração de material particulado, método automático,
configuração instrumental e comparador gravimétrico. Uma comparação poderá produzir
mais de uma estimativa, sem que essas estimativas sejam consideradas independentes.

\section{Coleta de dados}

Um formulário padronizado foi testado em uma amostra e revisado antes da extração dos
documentos disponíveis. A extração produzida com assistência de IA será considerada a
primeira rodada. O pesquisador responsável verificará os campos diretamente contra os
textos e suplementos em uma segunda rodada, priorizando os dados quantitativos, a
identificação dos equipamentos, a colocalização, a coincidência temporal e o número de
observações. A forma de conferência e as correções serão declaradas no registro e no
relatório final.

\subsection{Variáveis bibliográficas e contextuais}

\begin{itemize}
  \item autores, ano, país e fonte de publicação;
  \item coordenadas ou descrição do local, quando disponíveis;
  \item ambiente urbano, industrial, rural, de tráfego ou de fundo;
  \item período, duração, estações do ano e condições meteorológicas;
  \item fração de material particulado e faixa de concentração;
  \item fontes predominantes ou composição do aerossol, quando informadas.
\end{itemize}

\subsection{Características dos métodos}

\begin{itemize}
  \item fabricante, modelo, princípio e situação regulatória do equipamento automático;
  \item método gravimétrico, amostrador, entrada seletiva, vazão, filtro e duração;
  \item distância horizontal, diferença de altura e descrição da colocalização;
  \item resolução original e regra de agregação do equipamento automático;
  \item aquecimento, secagem, correções, calibração e manutenção;
  \item critérios de validade, completude, exclusão de dados e controle de qualidade;
  \item número de observações pareadas e tratamento de dados ausentes.
\end{itemize}

\subsection{Resultados}

Serão extraídos, sempre que disponíveis:

\begin{itemize}
  \item médias, desvios-padrão, medianas e intervalos dos dois métodos;
  \item diferença média e viés relativo;
  \item razão entre as concentrações;
  \item inclinação, intercepto e incerteza da regressão;
  \item limites de concordância;
  \item RMSE, MAE e outras medidas de erro;
  \item coeficientes de correlação e determinação;
  \item precisão, repetibilidade e incerteza;
  \item fatores de correção aplicados;
  \item atendimento ou não a critérios formais de equivalência.
\end{itemize}

Resultados obtidos diretamente de tabelas terão prioridade. Extrações de figuras,
conversões de unidades, cálculos derivados e respostas de autores serão identificados
separadamente.

\section{Análise bibliométrica descritiva}

O corpus bibliométrico compreenderá todos os documentos que atenderem aos critérios de
elegibilidade ao final das buscas de citações. Na data desta emenda, o corpus contém 40
documentos. A unidade bibliométrica será o documento, enquanto a síntese de evidências
utilizará o estudo, a família de estudo ou a comparação elegível, conforme a análise.

Serão descritos, conforme a completude dos metadados, ano de publicação, fonte,
autores, países ou locais estudados, palavras-chave, frações de material particulado,
tecnologias automáticas e tipos de comparador gravimétrico. Poderão ser apresentados
gráficos de evolução temporal, distribuição geográfica e frequência das tecnologias.
Redes de colaboração ou coocorrência serão exploradas somente se o tamanho e a
qualidade do corpus produzirem resultados interpretáveis.

Registros excluídos e publicações não recuperadas não integrarão a bibliometria
principal, embora permaneçam contabilizados no fluxo da revisão. Indicadores de
citação, produtividade autoral ou centralidade de periódicos não serão usados para
selecionar, excluir ou ponderar documentos. Não será aplicado refinamento
multicritério para reduzir o conjunto elegível. O processamento será realizado com
rotinas reproduzíveis, preferencialmente no ambiente R com o pacote
\texttt{bibliometrix}, e o código será disponibilizado.

\section{Desfechos e medidas de comparabilidade}

O desfecho principal será a direção e a magnitude da diferença entre o método
automático e o gravimétrico. Quando os dados permitirem, o viés relativo será expresso
por:

\begin{equation}
  B_{\%}=100\left(\frac{C_{\mathrm{auto}}-C_{\mathrm{grav}}}
  {C_{\mathrm{grav}}}\right),
\end{equation}

em que $C_{\mathrm{auto}}$ representa a concentração obtida pelo método automático e
$C_{\mathrm{grav}}$ a concentração gravimétrica correspondente. Valores positivos
indicam superestimação pelo método automático e valores negativos indicam
subestimação.

As medidas serão hierarquizadas da seguinte forma:

\begin{enumerate}
  \item viés relativo ou diferença média baseada em observações pareadas;
  \item razão entre métodos e limites de concordância;
  \item inclinação e intercepto de regressões de comparação;
  \item medidas de erro e precisão;
  \item correlação e $R^2$, tratados somente como medidas de associação complementar;
  \item classificação de equivalência conforme critério declarado pelo estudo.
\end{enumerate}

Não serão combinados em uma única estimativa resultados que representem propriedades
estatísticas distintas, como correlação, inclinação da regressão e diferença média.

\section{Caracterização e avaliação metodológica}

Será utilizada uma avaliação por domínios, específica para estudos de comparação de
métodos ambientais. A versão emendada compreenderá:

\begin{enumerate}
  \item seleção e descrição do local e do período;
  \item colocalização e coincidência temporal;
  \item adequação e rastreabilidade do método gravimétrico;
  \item calibração, manutenção e controle de qualidade do método automático;
  \item completude, exclusões e tratamento dos dados;
  \item cobertura da faixa de concentração e das condições ambientais;
  \item adequação das análises de concordância e comparabilidade;
  \item transparência e completude do relato;
  \item risco de relato seletivo e participação de financiadores ou fabricantes.
\end{enumerate}

Cada domínio será classificado como baixo risco, algumas preocupações, alto risco ou
informação insuficiente, acompanhado de justificativa textual. Não será produzida uma
pontuação total pela simples soma dos domínios. Os termos serão utilizados como
sinalizações metodológicas específicas desta revisão, e não como aplicação de um
instrumento externo validado. O pesquisador responsável avaliará as comparações e
repetirá a conferência em uma segunda rodada, sem consultar inicialmente o julgamento
anterior quando isso for operacionalmente possível. O teste-piloto assistido por IA
será tratado como primeira aplicação e não como avaliação independente.

A unidade operacional de julgamento será a comparação elegível, e não necessariamente
a publicação inteira. Informações comuns ao estudo ou ao relatório poderão ser
reutilizadas, mas diferenças entre locais, períodos, instrumentos, configurações ou
subconjuntos receberão julgamentos próprios. Será mantida também uma síntese por
família de estudo para evitar que publicações ou comparações múltiplas aumentem
artificialmente o peso da evidência.

As categorias de julgamento serão operacionalizadas da seguinte forma:

\begin{itemize}
  \item \textbf{baixo risco}: informação suficiente e ausência de limitação com efeito
  material provável sobre a comparação;
  \item \textbf{algumas preocupações}: omissão ou limitação plausivelmente relevante,
  sem evidência de comprometimento grave;
  \item \textbf{alto risco}: limitação com probabilidade substancial de alterar a
  direção ou a magnitude da comparação;
  \item \textbf{informação insuficiente}: relato que não permite julgamento, sem que
  isso seja convertido automaticamente em baixo ou alto risco.
\end{itemize}

Para cada domínio foram definidas regras específicas de julgamento, registradas no
instrumento de avaliação. Os domínios de seleção do local e cobertura ambiental serão
interpretados como ameaças à validade da comparação somente quando a seleção puder
alterar a direção ou a magnitude do resultado; limitações exclusivas de generalização
serão registradas separadamente como aplicabilidade.

Uma sinalização global, usada apenas nas análises de sensibilidade, será classificada
como alto risco quando houver alto risco em colocalização e coincidência temporal,
adequação do comparador gravimétrico, controle do método automático, completude dos
dados ou análise de comparabilidade, ou quando dois ou mais dos demais domínios forem
classificados como alto risco. Na ausência dessas condições, preocupações ou informação
insuficiente nos domínios 1 a 8 resultarão em sinalização global de algumas
preocupações.

A ausência isolada de declaração formal de financiamento ou conflitos de interesse,
frequente em publicações antigas, será mantida como informação insuficiente no domínio
9 e como marcador de transparência, mas não elevará automaticamente a sinalização
global. Participação de fabricante, seleção plausível de resultados ou influência do
financiador sobre desenho, análise ou relato continuarão contribuindo para a
sinalização global e para análises de sensibilidade. Baixo risco global exigirá baixo
risco nos domínios 1 a 8 e ausência de indício material de relato seletivo ou influência
indevida no domínio 9.

O teste-piloto foi aplicado a 40 documentos e, pelas regras anteriores, resultou em
quatro sinalizações de alto risco e 36 de algumas preocupações; em 30 documentos, a
informação sobre conflitos de interesse foi considerada insuficiente. Esses resultados
serão preservados como piloto. As sinalizações globais serão recalculadas no nível da
comparação depois da segunda conferência interna da versão emendada do instrumento.
Alterações serão registradas sem apagar as justificativas da primeira rodada. A
ausência de avaliação humana independente será considerada limitação do processo.

\pendente{Concluir a segunda conferência interna do instrumento emendado, registrar as
alterações em relação ao piloto e ratificar as regras antes da síntese definitiva.}

\section{Plano de síntese}

\subsection{Preparação e agrupamento}

Os estudos serão inicialmente agrupados por:

\begin{enumerate}
  \item fração: \MPdez{} ou \MPdoiscinco{};
  \item tecnologia automática: BAM, TEOM, TEOM-FDMS, óptica ou outra;
  \item natureza do método gravimétrico e situação de referência;
  \item artigo revisado por pares ou relatório técnico oficial;
  \item escala temporal da comparação;
  \item tipo de ambiente monitorado.
\end{enumerate}

Conversões de unidade serão documentadas. Quando um estudo apresentar múltiplos
equipamentos, locais, estações ou períodos, sua estrutura de dependência será
preservada; estimativas do mesmo estudo não serão tratadas automaticamente como
evidências independentes.

Todas as estimativas elegíveis serão mantidas na matriz longa. Para sínteses
quantitativas, relatórios da mesma base de dados serão reunidos sob uma única família
de estudo, utilizando-se o relatório mais completo e os demais como fontes
complementares. Se houver estimativas múltiplas dentro do mesmo estrato, será preferida
modelagem multinível ou estimação robusta da variância quando a quantidade de estudos
independentes e a estrutura dos dados permitirem. Caso isso não seja viável, será
selecionada uma estimativa por estudo e estrato, priorizando, nesta ordem: comparação
direta e estritamente colocalizada; coincidência temporal exata; período completo com
maior completude; estimativa acompanhada de incerteza e número de pares; e resultado
obtido de tabela ou texto em vez de digitalização de figura. Escolhas alternativas
plausíveis serão examinadas em análise de sensibilidade.

\subsection{Síntese descritiva}

Todos os estudos elegíveis integrarão tabelas de características e resultados. Serão
apresentados mapas ou gráficos de distribuição geográfica, cronologia das publicações,
tecnologias avaliadas, duração, quantidade de pares, faixa de concentração e direção
dos vieses. Os dados individuais serão exibidos sempre que isso favorecer a
interpretação e a rastreabilidade. Os resultados bibliométricos serão apresentados
separadamente dos resultados de comparabilidade e não determinarão o peso atribuído às
evidências.

\subsection{Síntese quantitativa}

Uma meta-análise será considerada somente quando houver pelo menos três estudos
independentes com fração, tecnologia, comparador e medida de resultado suficientemente
compatíveis. O agrupamento será realizado com modelo de efeitos aleatórios. Serão
informados estimativa combinada, intervalo de confiança, $\tau^2$, $I^2$ e, quando
viável, intervalo de predição.

Razões poderão ser analisadas em escala logarítmica. Diferenças absolutas poderão ser
agrupadas quando as unidades e os períodos forem compatíveis. Vieses percentuais
calculados por definições diferentes não serão combinados sem harmonização
demonstrável. Regressões, correlações e limites de concordância serão sintetizados em
estratos próprios.

Quando houver múltiplas estimativas dependentes, será aplicado o procedimento definido
na preparação e no agrupamento. A decisão analítica para cada estrato será registrada
antes do cálculo de qualquer estimativa combinada.

\subsection{Heterogeneidade e fatores condicionantes}

Se houver quantidade suficiente de estudos, serão examinados:

\begin{itemize}
  \item tecnologia e modelo do equipamento;
  \item \MPdez{} versus \MPdoiscinco{};
  \item umidade, temperatura e estação do ano;
  \item aquecimento, secagem e fatores de correção;
  \item tipo de ambiente e fontes predominantes;
  \item faixa de concentração;
  \item duração, número de pares e completude;
  \item situação regulatória do método automático e do comparador;
  \item sinalização das limitações metodológicas.
\end{itemize}

Subgrupos e meta-regressões somente serão utilizados quando houver base conceitual,
quantidade de estudos e distribuição de dados suficientes. Resultados exploratórios
serão identificados como tais.

\subsection{Análises de sensibilidade}

Estão previstas análises que:

\begin{itemize}
  \item excluam comparações com sinalização global de alto risco metodológico;
  \item incluam apenas artigos revisados por pares;
  \item incluam apenas métodos com reconhecimento regulatório declarado;
  \item incluam apenas comparações com colocalização estrita e demonstrada;
  \item excluam estudos sem agregação temporal exatamente coincidente;
  \item excluam resultados calculados a partir de figuras ou dados incompletos;
  \item avaliem o efeito de estudos com pequena quantidade de pares;
  \item excluam avaliações conduzidas por fabricantes sem salvaguardas independentes;
  \item comparem escolhas alternativas entre estimativas dependentes do mesmo estudo.
\end{itemize}

\subsection{Síntese sem meta-análise}

Quando o agrupamento quantitativo não for apropriado, será utilizada síntese
estruturada, com regras explícitas para agrupamento, métrica padronizada, ordenação,
apresentação gráfica e interpretação. A contagem simples de resultados
``significativos'' não será utilizada como medida de evidência. A direção, magnitude,
precisão e limitações metodológicas dos resultados permanecerão visíveis.

\subsection{Vieses de relato e confiança na evidência}

Assimetria de funil e testes associados somente serão avaliados quando houver pelo
menos dez estimativas comparáveis e quando as premissas desses procedimentos forem
razoáveis. A confiança no conjunto de evidências será discutida por resultado com base
nas limitações metodológicas, consistência, precisão, aplicabilidade direta e
possibilidade de resultados ausentes. A condução por um único pesquisador e as 17
publicações não recuperadas serão consideradas nessa discussão. Não será atribuída
automaticamente uma classificação clínica a estudos de desempenho instrumental.

\section{Disponibilidade, ética e disseminação}

A revisão utilizará documentos públicos e não envolverá participantes humanos nem
coleta de dados identificáveis. Serão disponibilizados, conforme as restrições das
bases e dos editores:

\begin{itemize}
  \item protocolo e histórico de emendas;
  \item estratégias completas de busca;
  \item formulários de seleção e extração;
  \item lista de textos completos excluídos e motivos;
  \item lista das publicações não recuperadas;
  \item dados extraídos, decisões de avaliação metodológica e registro das correções
  realizadas na segunda conferência interna;
  \item matriz bibliométrica e código das análises, tabelas e figuras;
  \item checklist PRISMA preenchido.
\end{itemize}

Os resultados serão submetidos como artigo científico e poderão subsidiar a
dissertação e o procedimento técnico relacionado ao planejamento de campanhas de
monitoramento. A revisão não estabelecerá equivalência regulatória por conta própria;
ela sintetizará as evidências e os critérios utilizados nas fontes incluídas.

\section{Decisões necessárias antes do registro e da síntese definitiva}

\begin{enumerate}
  \item decidir definitivamente sobre relatórios técnicos oficiais;
  \item fechar a lista de portais institucionais, caso haja literatura cinzenta;
  \item realizar as buscas retrospectiva e prospectiva de citações dos cinco
  estudos incluídos na busca adicional da PubMed e dos demais estudos elegíveis;
  \item realizar a segunda conferência interna da seleção, da extração e da avaliação
  metodológica, preservando o registro das correções;
  \item estruturar a matriz longa com identificadores de família, relatório,
  comparação e estimativa antes da síntese quantitativa;
  \item preparar e validar a matriz bibliométrica dos documentos incluídos;
  \item escolher o periódico-alvo ou ao menos seus limites gerais de extensão;
  \item adicionar ao Zotero as referências metodológicas de PRISMA-P, PRISMA-S,
  SWiM, CEE e análise bibliométrica e atualizar a exportação bibliográfica;
  \item aprovar a versão 1.0 e registrá-la com declaração explícita do estágio já
  alcançado, antes da síntese definitiva.
\end{enumerate}

\printbibliography[title={Referências disponíveis na biblioteca do projeto}]

\end{document}